\documentclass[11pt]{article}

\usepackage{amsmath,amssymb,graphicx}
\usepackage[margin=1in]{geometry}
\usepackage{booktabs}

\title{Mutual Registration and Constraint-Coupled Interaction Probe on a Kernel-Induced Cochain Operator Architecture}
\author{Tomislav Rupi\'c \\ Independent Researcher}
\date{March 2026}

\begin{document}

\maketitle

\begin{abstract}
Stage 21 opens the first post-reinforcement branch in the Phase II HAOS-IIP program. Stages 19A--19D and the Stage 20 boundary report established that weak substrate-law reinforcement, across local, mesoscopic, directional, and mixed symmetry classes, is insufficient to generate persistence in the current pre-geometric architecture. Stage 21 therefore asks a different question. Instead of strengthening the substrate, can persistence emerge when interaction itself becomes conditional on mutual registration between packet structures? The tested law class introduces bounded, memoryless interaction gating of the form $F_{AB}=\varepsilon R_{\mathrm{align}}F_{\mathrm{base}}$, where the registration functional is evaluated from instantaneous relational compatibility and no substrate memory is introduced. Three minimal registration classes are tested on the canonical representative triad consisting of a clustered composite anchor, a phase-ordered symmetric triad, and a counter-propagating corridor: phase-coherence registration, spatial-configuration registration, and composite-topology registration. Each class is evaluated under null, very weak, and weak interaction eligibility, giving a 27-run base pilot at resolution $n=12$. The result is uniformly negative at the persistence level. No composite lifetime gain appears, no binding persistence gain appears, no coarse basin persistence gain appears, no corridor dwell or locking gain appears, no new stabilized ordering class appears, and no refinement promotion is earned. The negative is still structurally informative because the registration channels are active. Phase-coherence registration reaches mean strengths up to $0.1829$ and effective interaction-energy fractions up to $3.3\times10^{-3}$, but duty cycle remains zero. Spatial-configuration registration reaches mean strengths up to $0.5519$, duty cycles up to $1.0$, and effective interaction-energy fractions up to $1.34\times10^{-2}$, yet produces no persistence gain. Composite-topology registration reaches mean strengths up to $0.3119$, duty cycles up to $1.0$, and effective interaction-energy fractions up to $7.27\times10^{-3}$, also without persistence consequences. The bounded conclusion is that selective interaction eligibility, in the tested weak and memoryless registration class, is insufficient to create persistent capture or geometry-like stabilization on this architecture.
\end{abstract}

\section{From Reinforcement to Registration}

Stage 20 closed a large region of Phase II mechanism space. Local scalar substrate reinforcement failed. Mesoscopic basin-residence reinforcement failed. Transport-corridor reinforcement failed. Flux-driven directional substrate response failed. Mixed scalar-vector reinforcement failed. In every case the branch was numerically active, deformation remained bounded and nonzero, and the persistence ladder stayed flat. The resulting boundary statement was clear: weak substrate-law reinforcement, across local, mesoscopic, directional, and mixed symmetry classes, is insufficient to generate persistence in this architecture at the tested couplings.

Stage 21 is the next coherent move because it leaves reinforcement space entirely. The substrate is not modified. Instead, the interaction law itself is gated by a bounded registration functional. The tested hypothesis is narrow. Perhaps packet structures do not need a stronger ambient medium; perhaps they need a selective interaction eligibility condition. If persistence requires relational compatibility rather than always-on weak force, then interaction should become active only when packet states enter a compatible region of configuration--phase space.

\section{Stage 21 Law Class}

The Stage 21 pilot uses a bounded interaction gate of the form
\[
F_{AB}(t)=\varepsilon\,R_{\mathrm{align}}(A,B,t)\,F_{\mathrm{base}}(A,B,t),
\]
with $\varepsilon\in\{0.01,0.02\}$ for the non-null runs and a null control at $\varepsilon=0$. The effective pairwise term is implemented as a registration-gated alignment drive projected onto the existing transverse sector. The operator, kernel support, normalization conventions, packet family, and promotion rules are all unchanged from the late Phase I and early Phase II pilots. Crucially, no cumulative substrate memory is introduced. The gate is instantaneous, bounded, and reversible.

Three registration classes are tested.
\begin{enumerate}
\item \textbf{Phase-coherence registration.} Eligibility depends on state alignment and velocity alignment. In executable form it behaves like a bounded resonance-compatibility test.
\item \textbf{Spatial-configuration registration.} Eligibility depends on bounded separation and overlap-compatible geometry. This is the closest Stage 21 analogue to structured encounter gating.
\item \textbf{Composite-topology registration.} Eligibility depends on motif-level compatibility, combining symmetry of pair distances, overlap, and alignment into a bounded motif score.
\end{enumerate}

All three are evaluated on the same canonical representative set inherited from the late pre-geometric atlas:
\begin{itemize}
\item the clustered composite anchor,
\item the phase-ordered symmetric triad,
\item the counter-propagating corridor.
\end{itemize}
The total base pilot is therefore $3$ registration classes $\times$ $3$ representatives $\times$ $3$ conditions $=27$ runs.

\section{Persistence-First Result}

The persistence-first read is uniform across the entire base matrix. Every one of the 27 runs remains in the label ``no registration effect.'' More importantly, the quantitative promotion metrics remain flat against class-matched null controls:
\begin{itemize}
\item composite lifetime delta remains $0.0000$,
\item binding persistence delta remains $0.0000$,
\item coarse basin persistence delta remains $0.0000$,
\item corridor dwell delta remains $0.0000$,
\item no new stabilized ordering class appears,
\item no run earns promotion to the $12\rightarrow24$ refinement ladder.
\end{itemize}

This is therefore a regime-level negative, not an ambiguous outcome. Stage 21 does not merely fail to produce a dramatic effect. It fails under the same persistence-first observables that have defined the Phase I and Phase II boundary logic from the start.

\section{Why the Negative Is Still Real}

The null does not come from a dead gate. Non-null runs carry bounded registration signal and bounded gated interaction energy in all three classes.

\subsection{Phase-Coherence Registration}

Phase-coherence registration remains the weakest of the three classes in its activation consequences. Mean registration strength reaches up to $0.1829$, with maximum registration strength reaching $0.5588$, and effective interaction-energy fraction reaching $3.32\times10^{-3}$. Yet the duty cycle remains $0.0$ across all phase-coherence runs. In operational terms, phase-style compatibility never opens a sustained interaction window. The gate flickers numerically, but not structurally.

\subsection{Spatial-Configuration Registration}

Spatial registration is much stronger in kind. Mean registration strength reaches $0.5519$, maximum registration strength reaches $0.6971$, effective interaction-energy fraction reaches $1.343\times10^{-2}$, and the duty cycle reaches $1.0$ for the strongest clustered-anchor case. The encounter dwell mean also becomes long on the base horizon, reaching $1.4764$ in the strongest runs. So this branch clearly opens stable eligibility windows. Nevertheless, none of that converts into a persistence gain. The clustered composite anchor, the symmetric triad, and the corridor control all remain regime-identical to null at the promoted observables.

\subsection{Composite-Topology Registration}

Composite-topology registration occupies an intermediate position. Mean registration strength reaches $0.3119$, the duty cycle also reaches $1.0$ in the strongest non-null cases, effective interaction-energy fraction reaches $7.27\times10^{-3}$, and motif survival time reaches $1.4764$ on the pilot horizon. So pattern-level recognition also enters the dynamics. But again no persistence consequence follows. The gate can detect and weight compatible motifs without generating composite stabilization, basin capture, or corridor locking.

\section{Comparative Activation Hierarchy}

The three registration classes differ in how fully they activate, but not in their persistence outcome. The hierarchy is summarized below.

\begin{table}[h]
\centering
\begin{tabular}{@{}llll@{}}
\toprule
Registration class & Mean strength scale & Activation pattern & Persistence outcome \\
\midrule
Phase coherence & $\sim 10^{-1}$ & weak, duty cycle $0$ & negative \\
Spatial configuration & $\sim 5\times10^{-1}$ & strong, duty cycle up to $1$ & negative \\
Composite topology & $\sim 3\times10^{-1}$ & intermediate, duty cycle up to $1$ & negative \\
\bottomrule
\end{tabular}
\caption{Stage 21 registration hierarchy. Eligibility windows can open, but none of the tested classes produce a persistence gain, corridor locking gain, or new stabilized ordering class.}
\end{table}

This hierarchy matters because it blocks an easy under-activation excuse. The spatial and topology branches are not simply too weak to matter numerically. They are active enough to create nonzero duty cycles and nontrivial interaction-energy fractions, yet still fail to alter the regime-level observables.

\section{Structural Interpretation}

Stage 21 therefore sharpens the architectural picture in a new direction. The system now supports all of the following, under tested conditions:
\begin{itemize}
\item transient collective morphology,
\item weak relational ordering,
\item bounded substrate memory,
\item bounded registration-gated interaction windows.
\end{itemize}
What it still does not support is persistence. This matters because the failure is now no longer attributable only to a missing scalar well or an absent mesoscopic substrate response. Even when interaction becomes selective and conditional on relational compatibility, packet structures still do not enter a stable persistence-supporting regime.

The structural reading is therefore more specific than the Phase II reinforcement boundary. Weak substrate reinforcement was insufficient. Stage 21 now shows that weak selective interaction eligibility is also insufficient. The architecture can recognize compatibility without using that recognition to produce capture.

\section{Interpretation Boundary}

Stage 21 does not support claims about geometry, observation, awareness, or a fully emergent interaction law. It supports only one bounded conclusion: in the tested weak and memoryless registration class, selective interaction gating does not create persistence. The negative therefore applies to a specific and scientifically meaningful region of model space:
\begin{itemize}
\item weak registration-gated pairwise coupling,
\item no substrate memory,
\item no topology change,
\item no new operator sector,
\item persistence-first evaluation ladder unchanged.
\end{itemize}

This does not rule out every relational mechanism. It rules out the simplest registration class that could have supplied persistence without changing the architecture more deeply.

\section{Conclusion}

Stage 21 is the first direct test of selective interaction eligibility in the HAOS-IIP program. The question was whether persistence might require conditional interaction rather than stronger background force. The answer, in the tested class, is no.

Phase-coherence registration never opens a sustained eligibility window. Spatial-configuration registration opens strong and durable windows, yet produces no persistence gain. Composite-topology registration detects compatible motifs, yet also produces no persistence consequence. Across all 27 base runs, the persistence ladder remains exactly flat.

The resulting statement is therefore sharp:

\begin{quote}
Mutual registration and constraint-coupled interaction eligibility are insufficient, at the tested couplings, to improve composite lifetime, binding persistence, coarse basin persistence, or corridor locking in the current architecture.
\end{quote}

This leaves the program in a stronger, not weaker, position. Stage 21 closes another physically natural law class without resorting to parameter forcing. The next coherent move is no longer another weak gating variant. It is a genuinely new degree of freedom, constraint class, or operator sector.

\end{document}
